\section{问题分析}

四问虽然分别要求预测、限速和优化，但都依赖同一个核心映射：给定温区设定 $\bm S$ 和传送速度 $v$，计算焊接区域中心的炉温曲线 $T_c(t;\bm S,v)$。因此应先建立并标定这一正向模型，再在同一模型上逐问增加决策变量和目标函数。

总体建模与求解流程如图~\ref{fig:model-workflow} 所示。四个问题不再分别另建一套模型，而是先由附件数据标定统一的正向仿真器，再在同一物理模型上依次完成预测、限速与优化；最后用细网格复算和约束回代统一检验结果。

\begin{figure}[H]
  \centering
  \includegraphics[width=0.98\textwidth]{paper/model_workflow.pdf}
  \caption{回焊炉炉温曲线建模与优化的总体工作流程}
  \label{fig:model-workflow}
\end{figure}

\subsection{问题一：由实验曲线反推传热参数}

温区设定值描述的是加热装置的控制目标，不是移动焊点的瞬时温度。建模需先回答两个问题：焊点沿炉膛移动时接触怎样的环境温度，环境热量又怎样传到焊接区域中心。为此将模型分成“空间温度场”和“板内传热”两层。前者处理炉口、温区和间隙，后者处理热惯性。附件数据用于辨识换热参数；完成辨识后把参数冻结，再代入问题一的新工况进行预测。

\subsection{问题二：寻找一维可行域的右端点}

温区设定固定后，只有速度 $v$ 需要确定。每给定一个速度，正向模型可产生一条炉温曲线，进而算出峰值、斜率和特征时长。因此问题二可转化为一维约束极值问题。先扫描速度区间以免漏掉不连续的可行段，再对发生符号变化的约束裕量求根，便能同时说明“最大值是多少”和“是哪一条约束限制了速度”。

\subsection{问题三：在可行域内压缩热暴露}

第三问同时调节四组温度和速度，目标又由阈值穿越和峰值位置共同决定，难以写出稳定的解析梯度。适合采用全局搜索找到可行盆地，再用无导数局部算法精修。由于最优解很可能落在工艺边界上，粗时间步造成的峰值误差会直接改变可行性，故搜索结果必须在更细网格上重新排序和复核。

\subsection{问题四：先给出前沿，再说明选择偏好}

“面积小”和“左右对称”是两个不同目标，不宜在没有说明量纲与权重时直接相加。先构造能够同时反映持续时间差和形状差的镜像误差，再用 $\varepsilon$-约束法生成 Pareto 解集。这样既能观察两目标的交换关系，也能在需要唯一工艺时依据明确偏好选择面积优先、均衡或对称优先方案。
